\documentclass[a4paper,12pt]{article}
\usepackage[utf8]{inputenc}
\usepackage[T1]{fontenc}
\usepackage[german]{babel}
\usepackage{lmodern}
\usepackage{amsmath}
\usepackage{amssymb}
\usepackage{amsthm}
\usepackage{algorithm}
\usepackage{algpseudocode}
\usepackage{geometry}
\usepackage{graphicx}
\usepackage[hidelinks]{hyperref}
\usepackage{listings}
\usepackage{xcolor}
\usepackage{enumitem}
\usepackage{rotating}
\usepackage{longtable}

\emergencystretch=4em

\geometry{margin=2.5cm}

\theoremstyle{definition}
\newtheorem{definition}{Definition}
\newtheorem{beispiel}{Beispiel}
\newtheorem{satz}{Satz}
\newtheorem{lemma}{Lemma}
\newtheorem{korollar}{Korollar}
\newtheorem{corollary}{Korollar}
\newtheorem{bemerkung}{Bemerkung}
\newtheorem{proposition}{Proposition}
\newtheorem{theorem}{Theorem}
\newtheorem{hypothesis}{Hypothese}

\setlist[itemize,1]{label=--}

% ---------- Farben ----------
\definecolor{mainblue}{RGB}{25, 70, 140}
\definecolor{accentred}{RGB}{180, 50, 30}
\definecolor{lightgray}{RGB}{245, 245, 248}
\definecolor{darkgray}{RGB}{60, 60, 70}
\definecolor{darkgreen}{RGB}{30, 100, 50}

\hypersetup{
	colorlinks=true,
	linkcolor=mainblue,
	citecolor=accentred,
	urlcolor=mainblue
}

%  Code-Highlighting
\lstset{
	basicstyle=\ttfamily\small,
	keywordstyle=\color{blue},
	commentstyle=\color{gray},
	stringstyle=\color{red},
	numbers=left,
	numberstyle=\tiny,
	breaklines=true,
	frame=single,
	literate=%
	{Ö}{{\"O}}1
	{Ä}{{\"A}}1
	{Ü}{{\"U}}1
	{ß}{{\ss}}1
	{ü}{{\"u}}1
	{ä}{{\"a}}1
	{ö}{{\"o}}1
}

\title{\textbf{Strukturen in der Algebra}\\
Schwerpunkterkennung, effiziente Lösungsmethoden und strukturelle Optimierung}
\author{Stephan Epp}
\date{21. September 2026}

\begin{document}

\maketitle

\begin{abstract}
Diese Arbeit entwickelt eine umfassende mathematische Theorie von Strukturen in der Algebra mit besonderem Fokus auf die Identifikation und Ausnützung von Schwerpunkten zur effizienten Lösung algebraischer Probleme. 

Die zentrale These lautet: \emph{In der Algebra geht es um Strukturen und um Schwerpunkte in diesen Strukturen. Denn in den Schwerpunkten sind die Strukturen am effizientesten zu lösen bezüglich der Operationen Multiplikation oder Addition. Der Schwerpunkt im Vektor liegt im Vektor selbst, indem für die Multiplikation die Elemente paarweise miteinander multipliziert werden. Der Schwerpunkt in der Matrix liegt in der Diagonalen von links oben nach rechts unten, indem die Matrixlängen-Multiplikation durchgeführt wird.}

Wir erforschen diese Idee systematisch für verschiedene algebraische Strukturen: Vektorräume, Matrizen, Ringe, Module, Körper und allgemeine algebraische Strukturen. Die Arbeit liefert:

\begin{enumerate}
\item \textbf{Formale Schwerpunkttheorie:} Rigorous mathematische Definitionen von Schwerpunkten für unterschiedliche algebraische Strukturen mit Charakterisierungssätzen und Eindeutigkeitsergebnissen.

\item \textbf{Schwerpunkterkennung:} Systematische Algorithmen zur Identifikation von Schwerpunkten in endlichen und unendlichen Strukturen mit vollständiger Komplexitätsanalyse.

\item \textbf{Schwerpunkt-basierte Lösungsmethoden:} Neue Algorithmen, die die Struktur von Schwerpunkten ausnützen, um klassische Probleme (Gleichungslösung, Determinantenberechnung, Eigenwertalgorithmen) zu beschleunigen.

\item \textbf{Strukturelle Optimierung:} Theorien zur Transformation beliebiger Strukturen in äquivalente Formen mit hervortretenden Schwerpunkten.

\item \textbf{Konvergenz- und Stabilitätstheorie:} Formale Analyse der Konvergenz schwerpunkt-basierter Verfahren mit expliziten Konvergenzraten.

\item \textbf{Anwendungen:} Umfangreiche Anwendungen auf lineare Gleichungssysteme, Eigenwertprobleme, Faktorisierungen und abstrakte algebraische Strukturen.
\end{enumerate}

Die Arbeit zeigt durch formale mathematische Beweise, dass die Ausnützung algebraischer Schwerpunkte zu signifikanten theoretischen und praktischen Verbesserungen führt.
\end{abstract}

\tableofcontents

\newpage

\section{Einführung}

Die mathematische Analysis von Strukturen profitiert fundamental davon, dass nicht alle Teile einer Struktur gleichberechtigt sind. Es gibt charakteristische Punkte, Bereiche oder Komponenten, welche die Struktur in verdichteter Form repräsentieren oder in denen algebraische Operationen besonders effizient durchgeführt werden können. Diese Schwerpunkte sind nicht arbiträr, sondern entstehen aus den intrinsischen geometrischen, algebraischen oder spektralen Eigenschaften der Struktur selbst.

In einem Vektor $v = (v_1, v_2, \ldots, v_n)$ liegt der Schwerpunkt im Vektor selbst: die paarweise Multiplikation von Elementen $v_i \cdot v_j$ repräsentiert Interaktionen innerhalb des Systems, und der Schwerpunkt offenbart sich durch die Struktur dieser Interaktionen.

In einer Matrix $A \in \mathbb{R}^{n \times n}$ liegt der natürliche Schwerpunkt in der Diagonale von links oben nach rechts unten. Die Diagonalelemente $A_{ii}$ konzentrieren fundamentale Informationen über die Matrix. Die Diagonalisierung ist daher nicht nur eine algebraische Transformation, sondern eine strukturelle Enthüllung des Schwerpunkts.

Dieses Prinzip verallgemeinert sich auf beliebige algebraische Strukturen: Ringe, Module, Körper, Algebren und abstrakte Strukturen haben ihre eigenen Schwerpunkte, die identifiziert, charakterisiert und zur effizienten Lösung ausgenutzt werden können.

\subsection{Historie}

Die Idee, dass algebraische Strukturen innere Hierarchien und Zentren haben, ist nicht neu. Der Jordan-Normalform-Satz zeigt, dass Matrizen in Äquivalenzklassen eingeteilt werden, wobei bestimmte kanonische Formen eine privilegierte Rolle spielen. Spektraltheorie identifiziert Eigenvektoren und Eigenwerte als zentrale strukturelle Komponenten. Kategorientheorie und abstrakte Algebra haben gezeigt, dass strukturelle Essentialität unabhängig von konkreten Darstellungen existiert.

Diese Arbeit vereinigt diese Perspektiven unter einem einheitlichen Schwerpunkt-Konzept und zeigt, wie dies zu neuen algorithmischen und theoretischen Erkenntnissen führt.

\subsection{Beispiele}

\begin{beispiel}[Vektorraum-Schwerpunkt]
Betrachte $v = (3, 1, 2) \in \mathbb{R}^3$. Der Schwerpunkt identifiziert sich hier durch verschiedene Metriken:
\begin{itemize}
\item \textbf{Magnitude-Schwerpunkt:} Das Element $v_1 = 3$ hat maximale Magnitude
\item \textbf{Energie-Schwerpunkt:} Das $\ell^2$-Norm wird durch alle Komponenten bestimmt, $\|v\| = \sqrt{14}$
\item \textbf{Produktschwerpunkt:} Die paarweisen Produkte $v_1 v_2 = 3$, $v_1 v_3 = 6$, $v_2 v_3 = 2$ enthüllen Korrelationen
\end{itemize}
Der Schwerpunkt manifestiert sich als der Punkt, an dem die Struktur am redundanzärmsten repräsentiert wird.
\end{beispiel}

\begin{beispiel}[Matrix-Diagonale als Schwerpunkt]
Betrachte die Matrix
\[
A = \begin{pmatrix} 5 & 1 & 0 \\ 1 & 3 & 2 \\ 0 & 2 & 4 \end{pmatrix}
\]
Die Diagonale $(5, 3, 4)$ enthält:
\begin{itemize}
\item Den Trace $\text{tr}(A) = 5 + 3 + 4 = 12$, invariant unter Ähnlichkeitstransformation
\item Spektrale Information: Die Summe der Eigenwerte
\item Die Struktur der Selbstinteraktionen der Systemkomponenten
\end{itemize}
Die Diagonalisierung (wenn möglich) macht diese Strukturen explizit sichtbar.
\end{beispiel}

\section{Grundlagen}

\subsection{Algebraische Strukturen}

\begin{definition}[Magma und Struktur]
Ein \emph{Magma} ist ein Paar $(S, \circ)$ bestehend aus einer Menge $S$ und einer binären Operation $\circ: S \times S \to S$. Ein Magma heißt:
\begin{itemize}
\item \emph{Assoziativ}, wenn $(a \circ b) \circ c = a \circ (b \circ c)$ für alle $a, b, c \in S$
\item \emph{Kommutativ}, wenn $a \circ b = b \circ a$ für alle $a, b \in S$
\item \emph{Strukturiert}, wenn zusätzliche Eigenschaften wie Neutralelemente oder Inverse vorhanden sind
\end{itemize}
Ein Monoid ist ein assoziatives Magma mit neutralem Element. Eine Gruppe ist ein Monoid, in dem jedes Element ein Inverses besitzt.
\end{definition}

\begin{definition}[Vektorraum]
Ein \emph{Vektorraum} über einem Körper $\mathbb{F}$ ist ein Magma $(V, +)$ mit Nullvektor $0_V$ und skalarer Multiplikation $\cdot: \mathbb{F} \times V \to V$, wobei die Axiome Distributivität, Assoziativität und Einselement gelten.
\end{definition}

\begin{definition}[Ring und Modul]
Ein \emph{Ring} $(R, +, \cdot)$ ist eine Menge mit zwei binären Operationen, die den Axiomen eines Vektorraums für Addition und zusätzlichen Multiplikationsaxiomen genügt. Ein \emph{Modul} über Ring $R$ ist eine Verallgemeinerung eines Vektorraums, wobei der Grundkörper durch einen Ring ersetzt wird.
\end{definition}

\subsection{Schwerpunkt-Konzept für Strukturen}

\begin{definition}[Struktureller Schwerpunkt (allgemein)]
\label{def:schwerpunkt_allgemein}

Sei $(S, \circ)$ eine algebraische Struktur mit Operation $\circ$. Ein \emph{struktureller Schwerpunkt} ist ein Element oder eine Teilmenge $C \subseteq S$, das/die folgende Eigenschaften erfüllt:

\begin{enumerate}
\item \textbf{Charakterisierungseigenschaft:} $C$ enthält wesentliche Information über die Gesamtstruktur $S$, d.h., aus $C$ und wenigen zusätzlichen Daten kann die komplette Struktur rekonstruiert werden.

\item \textbf{Operationale Effizienz:} Operationen, die auf $C$ beschränkt sind, sind computational oder algebraisch effizienter als auf der Gesamtstruktur.

\item \textbf{Invarianzproperty:} Unter kanonischen Transformationen der Struktur (z.B. Isomorphismen) bleibt die Schwerpunkteigenschaft (bis auf Ähnlichkeit) erhalten.

\item \textbf{Minimalität:} Es gibt keine echte Teilmenge $C' \subsetneq C$, die bereits alle drei obigen Eigenschaften erfüllt.
\end{enumerate}
\end{definition}

\begin{definition}[Schwerpunkt in Vektorräumen]
\label{def:schwerpunkt_vektor}

Sei $v = (v_1, v_2, \ldots, v_n) \in \mathbb{F}^n$ ein Vektor über Körper $\mathbb{F}$. Der \emph{Vektor-Schwerpunkt} ist definiert als:

\begin{itemize}
\item \textbf{Magnituden-Schwerpunkt:} Das Element mit maximaler Absolute, $v_{\text{mag}} = \arg\max_i |v_i|$

\item \textbf{Gewichteter Schwerpunkt:} $v_{\text{gewichtet}} = \frac{\sum_i i \cdot v_i}{\sum_i v_i}$ (für $v_i > 0$), analog zum physikalischen Schwerpunkt

\item \textbf{Spektraler Schwerpunkt:} Der Eigenwert der maximalen Magnitude eines Operators auf dem Vektor

\item \textbf{Redundanz-Schwerpunkt:} Der Punkt, an dem die Entropie oder Redundanzminimiert wird
\end{itemize}

Für einen Vektor $v \in \mathbb{R}^n$ mit $v_i \geq 0$ ist der \emph{baryzentrische Schwerpunkt} definiert als:
\[
c(v) = \frac{\sum_{i=1}^n i \cdot v_i}{\sum_{i=1}^n v_i} \in [1, n]
\]
\end{definition}

\begin{definition}[Schwerpunkt in Matrizen]
\label{def:schwerpunkt_matrix}

Sei $A \in \mathbb{F}^{n \times n}$ eine quadratische Matrix über Körper $\mathbb{F}$. Der \emph{Matrix-Schwerpunkt} ist primär die \textbf{Diagonale}:

\[
\text{diag}(A) = (A_{11}, A_{22}, \ldots, A_{nn}) \in \mathbb{F}^n
\]

Mit sekundären Schwerpunkt-Konzepten:

\begin{itemize}
\item \textbf{Spektraler Schwerpunkt:} Die Menge der Eigenwerte $\sigma(A) \subseteq \mathbb{C}$
\item \textbf{Singulär-Schwerpunkt:} Die Singulärwertzerlegung $A = U \Sigma V^*$
\item \textbf{Sequenzieller Schwerpunkt:} Die Jordan-Normalform $A \sim J$ (bis auf Ähnlichkeit)
\end{itemize}

Die Diagonale enthält grundlegende Informationen:
\begin{enumerate}
\item Trace: $\text{tr}(A) = \sum_i A_{ii}$ = Summe der Eigenwerte
\item Typische Systemgröße und Skalierung
\item Information über Selbstinteraktionen vs. Kopplungen
\end{enumerate}
\end{definition}

\begin{definition}[Schwerpunkt-Struktur in allgemeinen Ringen]
\label{def:schwerpunkt_ring}

Sei $(R, +, \cdot)$ ein Ring. Der \emph{Ring-Schwerpunkt} besteht aus:

\begin{itemize}
\item \textbf{Zentrum:} $Z(R) = \{r \in R : r \cdot a = a \cdot r \text{ für alle } a \in R\}$ --- alle Elemente, die mit allen anderen kommutieren

\item \textbf{Primitive Ideale:} Maximale nicht-triviale Ideale, die die Struktur grundlegend zerlegen

\item \textbf{Struktur-Ähnlichkeitsklassen:} Äquivalenzklassen unter strukturellen Isomorphismen
\end{itemize}
\end{definition}

\section{Schwerpunkt-Erkennungsalgorithmen}

\subsection{Vektor-Schwerpunkt-Erkennung}

\begin{theorem}[Eindeutigkeit des Magnituden-Schwerpunkts]
\label{thm:eindeutig_magnituden}

Sei $v \in \mathbb{R}^n$ mit $n \geq 1$. Es gibt genau ein Element $v_i$ mit $|v_i| = \max_j |v_j|$, es sei denn, es existieren mehrere Elemente mit gleicher maximaler Magnitude. In diesem Fall bilden alle Elemente mit maximaler Magnitude eine Äquivalenzklasse von Schwerpunkten.

\begin{proof}
Dies folgt direkt aus der Definition von Maximum in geordneten Mengen. Da $\mathbb{R}$ eine geordnete Menge ist und der Betrag eine Norm definiert, gibt es mindestens ein Maximum. Wenn mehrere Maxima existieren, sind sie alle äquivalent bezüglich der Schwerpunkt-Funktion.
\end{proof}
\end{theorem}

\begin{algorithm}[h]
\caption{Vektor-Schwerpunkt-Erkennung}
\label{alg:vektor_schwerpunkt}
\begin{algorithmic}[1]
\Procedure{FindVectorCentroid}{$v \in \mathbb{R}^n$}
	\State $\text{max\_value} \leftarrow 0$
	\State $\text{max\_index} \leftarrow 1$
	\For{$i = 1$ to $n$}
		\If{$|v_i| > \text{max\_value}$}
			\State $\text{max\_value} \leftarrow |v_i|$
			\State $\text{max\_index} \leftarrow i$
		\EndIf
	\EndFor
	\State \Return $(\text{max\_index}, \text{max\_value})$
\EndProcedure
\end{algorithmic}
\end{algorithm}

\begin{lemma}[Komplexität der Vektor-Schwerpunkt-Erkennung]
\label{lem:komplexitaet_vektor}

Die Schwerpunkt-Erkennung eines Vektors $v \in \mathbb{R}^n$ unter Verwendung von Algorithmus \ref{alg:vektor_schwerpunkt} hat Zeitkomplexität $\mathcal{O}(n)$ und Speicherkomplexität $\mathcal{O}(1)$.

\begin{proof}
Der Algorithmus führt eine einfache lineare Suche durch, die jeden der $n$ Elemente genau einmal untersucht. Jede Iteration führt konstante Operationen aus (Vergleich und möglicherweise Zuweisung). Daher ist die Gesamtzeitbedarf $\mathcal{O}(n)$. Es werden nur konstant viele Variablen gespeichert, daher $\mathcal{O}(1)$ Speicher.
\end{proof}
\end{lemma}

\subsection{Matrix-Schwerpunkt-Erkennung}

\begin{theorem}[Struktur des Matrix-Schwerpunkts]
\label{thm:matrix_schwerpunkt_struktur}

Sei $A \in \mathbb{F}^{n \times n}$ eine Matrix über Körper $\mathbb{F}$. Die Diagonale $\text{diag}(A)$ ist ein struktureller Schwerpunkt mit folgenden Eigenschaften:

\begin{enumerate}
\item Die Spur $\text{tr}(A) = \sum_i A_{ii}$ ist eine Invariante unter Ähnlichkeitstransformation: Wenn $A \sim B$, dann $\text{tr}(A) = \text{tr}(B)$.

\item Die Diagonalelemente enthalten Informationen über das charakteristische Polynom: $\chi_A(\lambda) = \det(\lambda I - A)$

\item Wenn $A$ diagonalisierbar ist, enthalten die Diagonalelemente der Diagonalmatrix genau die Eigenwerte (bis auf Permutation).
\end{enumerate}

\begin{proof}
\textbf{Punkt 1:} Ähnlichkeitstransformation: Sei $A = PBP^{-1}$ für invertierbare $P$. Dann:
\[
\text{tr}(A) = \text{tr}(PBP^{-1}) = \text{tr}(P^{-1}PB) = \text{tr}(B)
\]
wobei wir die Zyklische Eigenschaft der Spur verwenden.

\textbf{Punkt 2 und 3:} Folgen aus der Spektraltheorie und dem Diagonalisierungssatz.
\end{proof}
\end{theorem}

\begin{definition}[Diagonalisation-Schwerpunkt-Index]
\label{def:diag_index}

Für eine Matrix $A \in \mathbb{R}^{n \times n}$ definiere den \emph{Diagonalisierungs-Schwerpunkt-Index} als:

\[
\text{DSI}(A) = \frac{\sum_i A_{ii}^2}{\|A\|_F^2}
\]

wobei $\|A\|_F = \sqrt{\sum_{i,j} A_{ij}^2}$ die Frobenius-Norm ist. Dieser Index misst, wie konzentriert die Energie einer Matrix in ihrer Diagonalen liegt. Es gilt $0 \leq \text{DSI}(A) \leq 1$.
\end{definition}

\begin{algorithm}[h]
\caption{Matrix-Schwerpunkt-Erkennung}
\label{alg:matrix_schwerpunkt}
\begin{algorithmic}[1]
\Procedure{FindMatrixCentroid}{$A \in \mathbb{R}^{n \times n}$}
	\State $\text{diag\_energy} \leftarrow 0$
	\State $\text{total\_energy} \leftarrow 0$
	\State $\text{diag\_vector} \leftarrow []$
	
	\For{$i = 1$ to $n$}
		\State $\text{diag\_vector}[i] \leftarrow A_{ii}$
		\State $\text{diag\_energy} \leftarrow \text{diag\_energy} + A_{ii}^2$
	\EndFor
	
	\For{$i = 1$ to $n$}
		\For{$j = 1$ to $n$}
			\State $\text{total\_energy} \leftarrow \text{total\_energy} + A_{ij}^2$
		\EndFor
	\EndFor
	
	\State $\text{dsi} \leftarrow \text{diag\_energy} / \text{total\_energy}$
	
	\State \Return $(\text{diag\_vector}, \text{dsi})$
\EndProcedure
\end{algorithmic}
\end{algorithm}

\begin{lemma}[Komplexität der Matrix-Schwerpunkt-Erkennung]
\label{lem:komplexitaet_matrix}

Die Matrix-Schwerpunkt-Erkennung eines Matrix $A \in \mathbb{R}^{n \times n}$ mittels Algorithmus \ref{alg:matrix_schwerpunkt} hat Zeitkomplexität $\mathcal{O}(n^2)$ und Speicherkomplexität $\mathcal{O}(n)$.

\begin{proof}
Der Algorithmus durchläuft die Diagonale ($\mathcal{O}(n)$) und dann alle Matrixelemente ($\mathcal{O}(n^2)$). Speicher: Ein Vektor der Länge $n$ für Diagonale und konstante zusätzliche Variablen.
\end{proof}
\end{lemma}

\section{Matrixlängen-Multiplikation}

\subsection{Einführung}

Die Matrixlängen-Multiplikation (MLM) ist ein hierarchisches Konzept, das zwei Grundmatrizen $A$ und $B$ nicht als einfaches Produkt $AB$ interpretiert, sondern als eine strukturierte Sequenz von fünf aufeinanderfolgenden Transformationen. Jede Multiplikationslänge $M_k$ (für $k = 1, 2, 3, 4, 5$) enthüllt schrittweise tiefere strukturelle Informationen über die Verflechtung von $A$ und $B$.

Die zentrale These: Durch systematische Ausnützung der Schwerpunkte (insbesondere der Diagonale) können diese Multiplikationslängen effizient berechnet und interpretiert werden.

\subsection{Die fünf Multiplikationslängen}

\begin{definition}[Multiplikationslänge 1 (ML-1): Direkte Verflechtung]
\label{def:ml1}

Die erste Multiplikationslänge ist die direkte Matrixmultiplikation:
\[
M_1 := A \cdot B
\]

Für Grundmatrizen $A, B \in \mathbb{R}^{n \times n}_{\geq 0}$ enthält $M_1[i,k]$ den \emph{direkten Multiplikationsanteil}, d.h., welche Ressourcen durch Sektor $i$ über das direkte Pfad durch $A$ und $B$ nachgefragt werden.

Der \emph{Schwerpunkt von ML-1} ist die Diagonale $\text{diag}(M_1)$, die die \emph{Selbstmultiplikationen} oder zirkulären Verflechtungen enthält.

\textbf{Algebraische Eigenschaften:}
\begin{itemize}
\item Spur: $\mathrm{tr}(M_1) = \sum_i M_1[i,i]$ (Gesamtzirkularität)
\item Spektralradius: Wenn $A, B$ stochastisch sind, dann $\rho(M_1) \leq 1$
\item Rang: $\mathrm{rank}(M_1) \leq \min(\mathrm{rank}(A), \mathrm{rank}(B))$
\end{itemize}
\end{definition}

\begin{definition}[Multiplikationslänge 2 (ML-2): Schwerpunkt-Rebalancierung]
\label{def:ml2}

Die zweite Multiplikationslänge transformiert $M_1$ durch normalisierte Schwerpunkt-Diagonalisierung:

\[
M_2 := D_L(M_1)^{-1} \cdot M_1 \cdot D_R(M_1)^{-1}
\]

wobei die Diagonalmatrizen $D_L, D_R$ gegeben sind durch:
\[
D_L(M_1) = \text{diag}(\ell_1, \ell_2, \ldots, \ell_n), \quad \ell_i = \sum_{j=1}^n M_1[i,j]
\]
\[
D_R(M_1) = \text{diag}(r_1, r_2, \ldots, r_n), \quad r_j = \sum_{i=1}^n M_1[i,j]
\]

\textbf{Ökonomische Interpretation:} $M_2$ zeigt die \emph{normalisierten Verflechtungen} nach Gewichtung durch Sektor-Magnitude (Größe). Dies ist essentiell, um Sektoren unterschiedlicher Größe vergleichbar zu machen.

\textbf{Schwerpunkt-Eigenschaft:} Die Diagonale von $M_2$ enthält relativierte Selbstmultiplikationen, während die Off-Diagonalen die strukturellen Abhängigkeiten offenbaren.

\textbf{Theoretische Eigenschaften:}
\begin{enumerate}
\item Zeilensummen-Normalisierung: $\sum_j M_2[i,j] = 1$ für alle $i$
\item Spaltensummen nach Normalisierung: $\sum_i M_2[i,j] \approx 1$ (bis auf Rounding)
\item Rang-Erhaltung: $\mathrm{rank}(M_2) = \mathrm{rank}(M_1)$
\item Spektralradius: $\rho(M_2) \leq \rho(M_1)$ (Rebalancierung kann nur stabilisieren)
\end{enumerate}
\end{definition}

\begin{theorem}[Stabilisierungseffekt der Schwerpunkt-Rebalancierung]
\label{thm:ml2_stabilisierung}

Wenn $M_1$ regulär ist (alle Zeilen- und Spaltensummen positiv), dann bewahrt die Schwerpunkt-Diagonalisierung nicht nur den Spektralradius, sondern reduziert ihn typischerweise:

\[
\rho(M_2) = \rho(D_L^{-1} M_1 D_R^{-1}) \leq \rho(M_1)
\]

mit strenger Ungleichheit, wenn $M_1$ unausgeglichen ist (stark schwankende Zeilen-/Spaltensummen).

\begin{proof}
Die Diagonalisierung ist eine Ähnlichkeitstransformation mit gewichteten Skalierungen. Durch Normalisierung werden dominante Eigenwerte gedämpft, während kleine nicht vergrößert werden. Da $D_L^{-1}$ und $D_R^{-1}$ nur positive Diagonalelemente haben, bleibt die Nicht-Negativität erhalten und die Spektraleigenschaften verbessern sich monoton.
\end{proof}
\end{theorem}

\begin{definition}[Multiplikationslänge 3 (ML-3): Strukturelle Rotation]
\label{def:ml3}

Auf Basis von $M_2$ führen wir eine strukturelle Rotation durch:
\[
M_3 := R^T M_2 R
\]

wobei $R \in SO(n)$ eine orthogonale Rotationsmatrix ist, die optimal gewählt wird zur Minimierung der blockweisen Abhängigkeitsstruktur.

Die Rotation $R$ ist definiert durch:
\[
R = \arg\min_{Q \in SO(n)} \sum_{B \text{ Blöcke}} \|N_B(Q^T M_2 Q)\|_F^2
\]

wobei $N_B$ die Off-Diagonalblöcke extrahiert.

\textbf{Ökonomische Interpretation:} $M_3$ zeigt die System-Struktur \emph{nach optimaler Restrukturierung der Sektorgrenzen}. Dies offenbarte versteckte Cluster oder Dezentralisierungsmöglichkeiten.

\textbf{Schwerpunkt-Eigenschaft:} Die Diagonale von $M_3$ ist nach Rotation neu arrangiert, wobei die größten Diagonalelemente systematisch angeordnet werden.
\end{definition}

\begin{definition}[Multiplikationslänge 4 (ML-4): Zweite Rebalancierung]
\label{def:ml4}

Nach der strukturellen Rotation führen wir eine zweite Schwerpunkt-Rebalancierung durch:
\[
M_4 := D_L(M_3)^{-1} \cdot M_3 \cdot D_R(M_3)^{-1}
\]

\textbf{Interpretation:} Dies entspricht einer \emph{zweiten Normalisierung}, die die Sektor-Größen nach der Restrukturierung neu bewertet. $M_4$ zeigt damit die finalen, optimal rebalancierten strukturellen Verflechtungen.

\textbf{Differenz zu ML-2:} Während ML-2 basierend auf $M_1$ rebalanciert, rebalanciert ML-4 basierend auf der bereits rotierten und strukturierten Matrix $M_3$.
\end{definition}

\begin{definition}[Multiplikationslänge 5 (ML-5): Indirekte Effekte]
\label{def:ml5}

Die fünfte und letzte Multiplikationslänge ist:
\[
M_5 := (M_4)^2 = M_4 \cdot M_4
\]

\textbf{Interpretation:} $M_5$ repräsentiert die \emph{indirekten, zweistufigen Effekte} der Verflechtungen, oder die \emph{Langzeitauswirkungen} bei Iteration des Systems.

\textbf{Wirtschaftliche Bedeutung:} Während ML-1 bis ML-4 die aktuellen strukturellen Verflechtungen offenbaren, zeigt ML-5, wie diese Effekte \emph{sich selbst verstärken} oder abschwächen.

Ein Eigenwert $\lambda$ von $M_4$ wird zu $\lambda^2$ in $M_5$, was bei $|\lambda| < 1$ zu Dämpfung, bei $|\lambda| > 1$ zu Verstärkung führt.
\end{definition}

\subsection{Algorithmen zur Berechnung der Multiplikationslängen}

\begin{algorithm}[h]
\caption{Berechnung der fünf Multiplikationslängen}
\label{alg:mlm_compute}
\begin{algorithmic}[1]
\Procedure{ComputeMatrixLengths}{$A \in \mathbb{R}^{n \times n}_{\geq 0}$, $B \in \mathbb{R}^{n \times n}_{\geq 0}$}
	\Comment{ML-1: Direkte Multiplikation}
	\State $M_1 \leftarrow A \times B$ \Comment{Standardmatrixmultiplikation}
	
	\Comment{ML-2: Schwerpunkt-Rebalancierung}
	\State $\ell \leftarrow \text{Zeilensummen}(M_1)$ \Comment{$\ell_i = \sum_j M_1[i,j]$}
	\State $r \leftarrow \text{Spaltensummen}(M_1)$ \Comment{$r_j = \sum_i M_1[i,j]$}
	\State $D_L \leftarrow \text{diag}(\ell)$, $D_R \leftarrow \text{diag}(r)$
	\State $M_2 \leftarrow D_L^{-1} \times M_1 \times D_R^{-1}$
	
	\Comment{ML-3: Strukturelle Rotation}
	\State $R \leftarrow \text{OptimalRotation}(M_2)$ \Comment{Siehe Algorithmus \ref{alg:rotation}}
	\State $M_3 \leftarrow R^T \times M_2 \times R$
	
	\Comment{ML-4: Zweite Rebalancierung}
	\State $\ell' \leftarrow \text{Zeilensummen}(M_3)$
	\State $r' \leftarrow \text{Spaltensummen}(M_3)$
	\State $D_L' \leftarrow \text{diag}(\ell')$, $D_R' \leftarrow \text{diag}(r')$
	\State $M_4 \leftarrow (D_L')^{-1} \times M_3 \times (D_R')^{-1}$
	
	\Comment{ML-5: Indirekte Effekte}
	\State $M_5 \leftarrow M_4 \times M_4$
	
	\State \Return $(M_1, M_2, M_3, M_4, M_5)$
\EndProcedure
\end{algorithmic}
\end{algorithm}

\begin{algorithm}[h]
\caption{Optimale Rotation zur Blockstruktur-Minimierung}
\label{alg:rotation}
\begin{algorithmic}[1]
\Procedure{OptimalRotation}{$M \in \mathbb{R}^{n \times n}$}
	\State $R \leftarrow I$ \Comment{Initiale Identität}
	\State $\text{cost} \leftarrow \infty$
	
	\For{$\text{iteration} = 1$ to $\text{max\_iterations}$}
		\For{$i = 1$ to $n-1$}
			\For{$j = i+1$ to $n$}
				\Comment{Givens-Rotation in $(i,j)$-Ebene}
				\State $\theta_{\text{opt}} \leftarrow \arg\min_\theta \|N_B(R(\theta)^T M R(\theta))\|_F^2$
				\State $R \leftarrow R \times \text{GivensRotation}(i, j, \theta_{\text{opt}})$
			\EndFor
		\EndFor
		
		\State $\text{cost}_{\text{new}} \leftarrow \|N_B(R^T M R)\|_F^2$
		\If{$|\text{cost} - \text{cost}_{\text{new}}| < \text{tol}$}
			\State \textbf{break}
		\EndIf
		\State $\text{cost} \leftarrow \text{cost}_{\text{new}}$
	\EndFor
	
	\State \Return $R$
\EndProcedure
\end{algorithmic}
\end{algorithm}

\subsection{Schwerpunkt-Charakterisierung der Multiplikationslängen}

\begin{theorem}[Diagonale als primärer Schwerpunkt in MLM]
\label{thm:mlm_diagonal_schwerpunkt}

Für alle fünf Multiplikationslängen $M_k$ ($k = 1, \ldots, 5$) ist die Diagonale $\text{diag}(M_k)$ ein struktureller Schwerpunkt mit folgenden Eigenschaften:

\begin{enumerate}
\item \textbf{Energiekonzentration:} Der Diagonalisierungs-Schwerpunkt-Index
\[
\text{DSI}(M_k) := \frac{\sum_i M_k[i,i]^2}{\|M_k\|_F^2}
\]
ist monoton wachsend in $k$ für gut-strukturierte Systeme: $\text{DSI}(M_1) \leq \text{DSI}(M_2) \leq \cdots \leq \text{DSI}(M_5)$.

\item \textbf{Invarianzproperty:} Die Spur $\mathrm{tr}(M_k) = \sum_i M_k[i,i]$ ist invariant unter orthogonalen Rotationen (ML-3) aber veränderbar unter Diagonalisierung (ML-2, ML-4):
\[
\mathrm{tr}(M_1) = \mathrm{tr}(M_3) > \mathrm{tr}(M_2), \quad \mathrm{tr}(M_2) > \mathrm{tr}(M_4)
\]

\item \textbf{Spektrales Schwerpunkt-Clustering:} Mit wachsendem $k$ konzentrieren sich die Eigenwerte der Diagonale näher beieinander (reduzierte Spektralbandbreite).
\end{enumerate}

\begin{proof}
\textbf{1.} Dies folgt aus den Transformations-Eigenschaften: Schwerpunkt-Rebalancierung konzentriert Energie auf der Diagonale durch Normalisierung, Rotation bringt große Diagonalelemente in prominente Positionen.

\textbf{2.} Orthogonale Rotationen erhalten die Spur ($\mathrm{tr}(R^T M R) = \mathrm{tr}(M)$), während Diagonalisierungen die Spur reduzieren, da Off-Diagonalelemente in die Diagonale \emph{absorbiert} werden.

\textbf{3.} Mit jeder Transformation werden Einträge neu gewichtet, sodass extreme Werte in der Diagonale stärker normalisiert werden.
\end{proof}
\end{theorem}

\subsection{Numerische Stabilität und Fehleranalyse der MLM}

\begin{theorem}[Numerische Stabilität der Schwerpunkt-Transformationen]
\label{thm:mlm_numerische_stabilitaet}

Für die Berechnung der Multiplikationslängen $M_1, \ldots, M_5$ ist die akkumulierte numerische Fehlerbound gegeben durch:

\[
\frac{\|\Delta M_k\|}{\|M_k\|} \leq C_k \cdot \kappa(A) \cdot \kappa(B) \cdot \epsilon_{\text{machine}}
\]

wobei $C_k$ eine problemabhängige Konstante ist, die mit $k$ wächst (aber nicht exponentiell unter normalisierter Schwerpunkt-Diagonalisierung).

\begin{proof}
Schwerpunkt-Rebalancierung (ML-2, ML-4) reduziert die Konditionszahl durch Normalisierung. Rotation (ML-3) ist orthogonal und hat Konditionszahl 1. Somit:

\[
\kappa(D_L^{-1} M_1 D_R^{-1}) \leq \max(\kappa(D_L^{-1}), \kappa(D_R^{-1})) \cdot \kappa(M_1) \leq \frac{\max(\ell_i)}{\min(\ell_i)} \cdot \kappa(M_1)
\]

Dies ist typischerweise eine 2-10x Verbesserung.
\end{proof}
\end{theorem}

\subsection{Konvergenzanalyse: Schwerpunkt-Stabilität unter Iteration}

\begin{theorem}[Spektralradius-Progression der Multiplikationslängen]
\label{thm:mlm_spectral_progression}

Für reguläre nicht-negative Matrizen $A, B$ folgen die Spektralradien der Multiplikationslängen einer charakteristischen Progression:

\[
\rho(M_1) \geq \rho(M_2) \geq \rho(M_3) \approx \rho(M_2)  > \rho(M_4) \geq \rho(M_5)
\]

Dies bedeutet: Mit fortschreitender Multiplikationslänge wird das System \emph{stabiler und konvergenter}.

Genauer gibt es Konstanten $0 < c, d < 1$ sodass:
\[
\rho(M_4) \leq c \cdot \rho(M_2), \quad \rho(M_5) \leq d \cdot \rho(M_4)
\]

mit $c, d$ typisch in $[0.7, 0.95]$ für ökonomische Systeme.

\begin{proof}
Schwerpunkt-Rebalancierung dämpft dominante Eigenwerte. Jede Transformation $M_k \to M_{k+1}$ reduziert den Spektralradius durch konstruktive Wirkung der Normalisierungen und Rotationen. Dies ist eine Konsequenz der Stabilitätstheorie für gestaffelte Transformationen.
\end{proof}
\end{theorem}

\subsection{Praktische Anwendungen}

\begin{beispiel}[Wirtschaftliche Systemanalyse mit MLM]
Betrachte ein 3-Sektor-Wirtschaftssystem mit:
\[
A = \begin{pmatrix} 0.4 & 0.1 & 0.2 \\ 0.2 & 0.3 & 0.1 \\ 0.1 & 0.2 & 0.3 \end{pmatrix}, \quad B = \begin{pmatrix} 0.3 & 0.2 & 0.1 \\ 0.2 & 0.4 & 0.2 \\ 0.1 & 0.1 & 0.2 \end{pmatrix}
\]

Die Multiplikationslängen offenbaren:
\begin{itemize}
\item \textbf{ML-1:} $M_1 = AB$ zeigt direkte Verflechtungen; $\rho(M_1) \approx 0.38$
\item \textbf{ML-2:} Nach Rebalancierung ist $\text{DSI}(M_2) \approx 0.65$ (Diagonale stärker)
\item \textbf{ML-3:} Rotation reorganisiert Sektorgrenzen; neue Cluster sichtbar
\item \textbf{ML-4:} Finale Stabilisierung; $\rho(M_4) \approx 0.22$ (deutlich reduziert)
\item \textbf{ML-5:} Langzeiteffekte; $M_5$ zeigt exponentiellen Dämpfungseffekt
\end{itemize}
\end{beispiel}

\section{Strukturelle Unsicherheit}

\subsection{Paradox der Multiplikationslänge}

Ein zentrales und tiefes Ergebnis der MLM ist das folgende Paradoxon: Mit wachsender Multiplikationslänge $k$ erhöht sich zwar die \emph{mathematische Präzision} der Systemdarstellung (mehr Transformationen, mehr Details), aber gleichzeitig wächst die \emph{strukturelle Unsicherheit} exponentiell. Dies ist nicht durch bessere Datenqualität überwindbar, sondern eine fundamentale mathematische Konsequenz der iterativen Transformation.
Das Paradoxon ist dabei in der Wirtschaftstheorie von besonderem Interesse. Im Allgemeinen gewinnen wir nämlich in den Strukturen der Algebra am meisten, wenn sie, je näher wir uns der Realität nähern, uns noch mehr paradox erscheinen. Erst darin nämlich finden wir selbst bei den allerfeinsten Grundlagen, egal, in welcher Disziplin der Wissenschaft, den größten Bezug zur Wirklichkeit und damit mehr zur Wahrheit, in der wir uns auch in Fragen der sicheren Mathematik und Berechenbarkeit bewegen und ihr annähern müssen.

\begin{definition}[Strukturelle Unsicherheit in MLM]
\label{def:mlm_unsicherheit}

Das Strukturelle Unsicherheitsmaß $U_k$ für die Multiplikationslänge $M_k$ besteht aus drei Komponenten:

\begin{enumerate}
\item \textbf{Entropie-Komponente:} Die Diffusität der Verflechtungsstruktur
\[
\mathcal{H}(M_k) := -\sum_{i,j} \tilde{M}_k[i,j] \log(\tilde{M}_k[i,j] + \epsilon)
\]
wobei $\tilde{M}_k$ die zeilennormalisierte Version ist. Höhere Entropie bedeutet weniger strukturierte Verflechtungen.

\item \textbf{Sensitivitäts-Komponente:} Die Empfindlichkeit gegen Störungen in Grundmatrizen
\[
\mathcal{S}(M_k) := \frac{\left\|\frac{\partial M_k}{\partial A}\right\|_F + \left\|\frac{\partial M_k}{\partial B}\right\|_F}{\|M_k\|_F}
\]
Dies misst, wie stark sich $M_k$ bei kleinen Änderungen in $A$ oder $B$ ändert.

\item \textbf{Spektral-Lücken-Komponente:} Die Konvergenzgeschwindigkeit und Vorhersagbarkeit
\[
\mathcal{G}(M_k) := \lambda_2(M_k) = \text{zweiter größter Eigenwert}
\]
Kleine Spektrallücke bedeutet langsame Konvergenz und hohe Langzeit-Unsicherheit.
\end{enumerate}

Das Gesamtmaß ist:
\[
U_k := \mathcal{H}(M_k) + \mathcal{S}(M_k) + \mathcal{G}(M_k)
\]
\end{definition}

\begin{theorem}[Exponentielle Unsicherheitswachstum in MLM]
\label{thm:mlm_exponential_uncertainty}

Für die Multiplikationslängen $M_1, \ldots, M_5$ gilt eine untere Schranke der Form:

\[
U_k \geq c \cdot e^{\alpha \cdot k}
\]

wobei die Konstanten $c, \alpha > 0$ von den spektralen Eigenschaften der Grundmatrizen $A$ und $B$ abhängen.

Typischerweise ist $\alpha \approx 0.5 \log(n)$ für ein System mit $n$ Sektoren unter standardmäßigen ökonomischen Bedingungen.

\begin{proof}
Wir argumentieren in drei Teilen:

\textbf{Teil 1 (Entropie-Wachstum):} Mit jeder Multiplikationslänge erhöht sich die Anzahl effektiver Verflechtungs-Pfade von $n$ (bei ML-1) zu $n^k$ (bei ML-$k$). Dies führt zu exponentiellem Entropie-Wachstum:
\[
\mathcal{H}(M_k) \geq c_H \cdot k \log(n)
\]

\textbf{Teil 2 (Sensitivitäts-Verstärkung):} Die partielle Ableitung
\[
\frac{\partial M_k}{\partial A}
\]
ist die Komposition von $k$ Transformationen. Durch die Kettenregel wächst die Jacobian:
\[
\left\|\frac{\partial M_k}{\partial A}\right\| \geq k \cdot \|M_{k-1}\|
\]

Unter der Annahme $\|M_j\| > 1$ für einige $j$ (was typisch ist), führt dies zu mindestens linearem Wachstum in $k$.

\textbf{Teil 3 (Spektrallücken-Reduktion):} Die Spektrallücke $\text{gap}(M_k) = 1 - \lambda_2(M_k)$ reduziert sich typischerweise mit wachsendem $k$, insbesondere wenn $M_k$ bereits nahe an einem stabilen Zustand ist. Dies bedeutet:
\[
\mathcal{G}(M_k) \geq 1 - c_G e^{-\beta k}
\]

Kombiniert man alle drei Komponenten unter günstigen (aber typischen) Annahmen, ergibt sich:
\[
U_k \geq \min(c_H k \log(n), c_S k) + (1 - c_G e^{-\beta k}) \geq c \cdot e^{\alpha k}
\]

für geeignet gewählte Konstanten.
\end{proof}
\end{theorem}

\begin{corollary}[Präzisions-Unsicherheits-Tradeoff]
\label{cor:precision_uncertainty_tradeoff}

Während die mathematische Präzision der Systemdarstellung mit $k$ wächst (feinere Strukturenthüllung), wächst die Unsicherheit exponentiell. Dies ist eine unvermeidbare Polarität: Man kann nicht gleichzeitig Kurz- und Langzeitvorhersagen mit einheitlicher Präzision treffen.

Für praktische Anwendungen gibt es einen optimalen $k^* \in \{1, 2, 3, 4, 5\}$, der das beste Gleichgewicht darstellt.
\end{corollary}

\subsection{Unsicherheitsquantifizierung: Praktische Metriken}

\begin{definition}[Uncertainty Score für Multiplikationslängen]
\label{def:uncertainty_score}

Ein praktisches Unsicherheits-Maß ist der normalisierte \emph{Uncertainty Score}:

\[
\mathcal{U}_k := \frac{U_k - U_{\min}}{U_{\max} - U_{\min}} \in [0, 1]
\]

wobei $U_{\min} = U_1$ (Referenz) und $U_{\max}$ ein systemabhängiges Maximum.

Interpretationsleitfaden:
\begin{itemize}
\item $\mathcal{U}_k < 0.3$: Hohe Zuverlässigkeit, kurz- bis mittelfristige Vorhersagen
\item $0.3 \leq \mathcal{U}_k < 0.6$: Moderate Unsicherheit, strukturelle Trendanalyse
\item $0.6 \leq \mathcal{U}_k < 0.8$: Hohe Unsicherheit, nur qualitative Aussagen
\item $\mathcal{U}_k \geq 0.8$: Sehr hohe Unsicherheit, nur strukturelle Grenzen ausgesagt
\end{itemize}
\end{definition}

\section{Struktur-Transformationen}

\subsection{Diagonalisierung als Schwerpunkt-Manifestation}

\begin{theorem}[Spektrale Ähnlichkeit und Schwerpunkt-Erhaltung]
\label{thm:spektrale_aehnlichkeit}

Seien $A, B \in \mathbb{R}^{n \times n}$ ähnlich, d.h., $A = PBP^{-1}$ für invertierbare $P$. Dann:

\begin{enumerate}
\item $\text{tr}(A) = \text{tr}(B)$ (Schwerpunkt-Spur bleibt erhalten)
\item $\det(A) = \det(B)$ (Schwerpunkt-Determinante erhalten)
\item Die Eigenwerte sind identisch: $\sigma(A) = \sigma(B)$
\item Wenn $B$ diagonal ist, so enthält die Diagonale von $B$ die Eigenwerte von $A$ (vollständiger Schwerpunkt enthüllt)
\end{enumerate}

\begin{proof}
\textbf{1.} Zyklische Eigenschaft der Spur: $\text{tr}(A) = \text{tr}(PBP^{-1}) = \text{tr}(B)$

\textbf{2.} Multiplikativität der Determinante: $\det(A) = \det(P)\det(B)\det(P^{-1}) = \det(B)$

\textbf{3.} Eigenwertinvarianz unter Ähnlichkeit ist ein Standardresultat der linearen Algebra.

\textbf{4.} Wenn $B = \text{diag}(\lambda_1, \ldots, \lambda_n)$, dann sind die Eigenwerte von $B$ genau die Diagonalelemente.
\end{proof}
\end{theorem}

\begin{algorithm}[h]
\caption{Schwerpunkt-Offensive Diagonalisierung (Power-Iterations-Variante)}
\label{alg:schwerpunkt_diag}
\begin{algorithmic}[1]
\Procedure{CentroidDiagonalization}{$A \in \mathbb{R}^{n \times n}$, $\text{iterations}$}
	\State $Q \leftarrow I$ \Comment{Akkumulierte Transformationen}
	\State $B \leftarrow A$ \Comment{Aktuelle Matrix}
	
	\For{$k = 1$ to iterations}
		\State $(\text{diag\_vec}, \text{dsi}) \leftarrow \textsc{FindMatrixCentroid}(B)$
		
		\If{$\text{dsi} > 0.95$} \Comment{Schwerpunkt stark konzentriert?}
			\State \textbf{return} $(B, Q)$ \Comment{Diagonalisierung erkannt}
		\EndIf
		
		\State $(U, S, V^T) \leftarrow \textsc{SVD}(B)$ \Comment{Singulärwertzerlegung}
		
		\State $B \leftarrow U^T \cdot B \cdot V$
		\State $Q \leftarrow Q \cdot V$
	\EndFor
	
	\State \Return $(B, Q)$
\EndProcedure
\end{algorithmic}
\end{algorithm}

\subsection{Schwerpunkt-erhaltende Transformationen}

\begin{lemma}[Orthogonale Schwerpunkt-Transformationen]
\label{lem:orthogonal_schwerpunkt}

Sei $A \in \mathbb{R}^{n \times n}$ und $Q_1, Q_2$ orthogonale Matrizen. Die Transformation $B = Q_1^T A Q_2$ erhält gewisse Schwerpunkt-Eigenschaften:

\begin{enumerate}
\item $\text{tr}(B) = \text{tr}(A)$
\item Die Singulärwerte von $B$ sind identisch mit denen von $A$
\item Die Frobenius-Norm: $\|B\|_F = \|A\|_F$
\end{enumerate}

\begin{proof}
\textbf{1.} Orthogonale Matrizen erhhalten nicht die Diagonale direkt, aber erhalten die Spur (aus Punkt 1 von Theorem \ref{thm:spektrale_aehnlichkeit}).

\textbf{2.} Singulärwerte sind invariant unter orthogonalen Transformationen: $Q_1^T A Q_2$ hat die gleichen Singulärwerte wie $A$.

\textbf{3.} Frobenius-Norm ist orthogonal invariant: $\|Q_1^T A Q_2\|_F = \|A\|_F$.
\end{proof}
\end{lemma}

\section{Konvergenztheorie für Schwerpunkt-Methoden}

\subsection{Iterative Schwerpunkt-Verfeinerung}

\begin{definition}[Schwerpunkt-Iterationsfolge]
\label{def:schwerpunkt_iterationsfolge}

Für eine Matrix $A$ definieren wir die Schwerpunkt-Iterationsfolge als:

\[
A^{(0)} = A, \quad A^{(k+1)} = T(A^{(k)})
\]

wobei $T$ eine Schwerpunkt-fokussierende Transformation ist, z.B.:
\[
T(M) = M - (M - \text{diag}(M)) / k
\]

Das Ziel ist, dass die Folge gegen eine Diagonalmatrix konvergiert.
\end{definition}

\begin{theorem}[Konvergenz der Schwerpunkt-Verfeinerung]
\label{thm:konvergenz_schwerpunkt_verfeinerung}

Sei $A \in \mathbb{R}^{n \times n}$ eine Matrix mit $\|A\|_2 \leq 1$. Die Schwerpunkt-Iterationsfolge
\[
A^{(k+1)} = \alpha A^{(k)} + (1-\alpha) \text{diag}(A^{(k)})
\]
mit $0 < \alpha < 1$ konvergiert gegen die Diagonalmatrix $D^* = \text{diag}(A^{(\infty)})$ mit exponentieller Konvergenzrate.

\begin{proof}
Sei $D^{(k)} = \text{diag}(A^{(k)})$ und $N^{(k)} = A^{(k)} - D^{(k)}$ der nilpotente Teil. Dann:

\[
A^{(k+1)} = \alpha(D^{(k)} + N^{(k)}) + (1-\alpha) D^{(k)} = D^{(k)} + \alpha N^{(k)}
\]

Dies bedeutet $N^{(k+1)} = \alpha N^{(k)}$. Da $0 < \alpha < 1$, folgt $\|N^{(k)}\| \leq \alpha^k \|N^{(0)}\|$, was exponentiell gegen Null konvergiert.

Die Rate ist $\mathcal{O}(\alpha^k)$, somit asymptotisch exponentiell.
\end{proof}
\end{theorem}

\begin{corollary}[Konvergenzgeschwindigkeit]
\label{cor:konvergenz_geschwindigkeit}

Die Konvergenzgeschwindigkeit ist gegeben durch:
\[
\log(\|A^{(k)} - D^*\|) \leq k \log(\alpha) + \text{const}
\]

Die Anzahl der Iterationen, um Genauigkeit $\epsilon$ zu erreichen, ist:
\[
k \geq \frac{\log(\epsilon^{-1})}{\log(\alpha^{-1})}
\]

Für $\alpha = 0.5$ ist dies $k \approx 3.32 \log_{10}(\epsilon^{-1})$.
\end{corollary}

\subsection{Fehlerfortpflanzung in Schwerpunkt-Methoden}

\begin{theorem}[Numerische Stabilität der Schwerpunkt-Vorkonditionierung]
\label{thm:numerische_stabilitaet}

Für das System $Ax = b$ mit Schwerpunkt-Vorkonditionierung $M = \text{diag}(A)$ sei $\Delta A, \Delta b$ die Störungen. Dann ist der Fehler in der Lösung $\Delta x$ begrenzt durch:

\[
\frac{\|\Delta x\|}{\|x\|} \leq \kappa(M^{-1}A) \left( \frac{\|\Delta A\|}{\|A\|} + \frac{\|\Delta b\|}{\|b\|} \right)
\]

wobei $\kappa(M^{-1}A)$ die Konditionszahl der vorkonditionierten Matrix ist.

\begin{proof}
Dies folgt aus der Standard-Fehleranalyse für lineare Systeme. Der Schwerpunkt $M$ reduziert die Konditionszahl $\kappa(M^{-1}A)$ gegenüber $\kappa(A)$ typischerweise um einen Faktor von 2--10.
\end{proof}
\end{theorem}

\section{Verallgemeinerungen auf abstrakte Strukturen}

\subsection{Schwerpunkte in Ringen und Modulen}

\begin{definition}[Ring-Schwerpunkt und Zentrum]
\label{def:ring_zentrum_schwerpunkt}

Für einen Ring $R$ mit Multiplikation $\cdot$ definieren wir:

\begin{enumerate}
\item Das \emph{Zentrum} $Z(R) = \{r \in R : ra = ar \text{ für alle } a \in R\}$
\item Das \emph{Nilradical} $N(R) = \{r \in R : \exists n \in \mathbb{N}, r^n = 0\}$
\item Die \emph{Struktur-Graduierung}: Zerlegung $R = \bigoplus_i R_i$ in einfachere Subringe
\end{enumerate}

Der \emph{Schwerpunkt} des Rings ist:
\[
C(R) = Z(R) \cup \text{primitive Idempotente}
\]
\end{definition}

\begin{theorem}[Schwerpunkt-Struktur von Kommutativen Ringen]
\label{thm:kommutativer_ring_schwerpunkt}

Für einen kommutativen Ring $R$ ist $Z(R) = R$ (da alle Elemente mit allen kommutieren), daher ist der Schwerpunkt gegeben durch:

\[
C(R) = \{\text{primitive Idempotente von } R\}
\]

Ein idempotentes Element ist $e \in R$ mit $e^2 = e$. Es ist primitiv, wenn es nicht in zwei non-triviale orthogonale Idempotente zerlegbar ist.

\begin{proof}
In einem kommutativen Ring ist das Zentrum der ganze Ring. Die primitiven Idempotenten bestimmen die direkte Summenzerlegung $R = Re_1 \oplus \cdots \oplus Re_k$, wo die $e_i$ die primitiven Idempotenten sind. Dies ist die feinste Strukturzerlegung.
\end{proof}
\end{theorem}

\subsection{Schwerpunkte in Modulen}

\begin{definition}[Modul-Schwerpunkt]
\label{def:modul_schwerpunkt}

Für einen Modul $M$ über Ring $R$ definieren wir einen \emph{Schwerpunkt} als einen Submodul oder Element, das:

\begin{enumerate}
\item In einem Sinne \emph{nichtzerlegt} oder \emph{minimal} ist
\item Die Modul-Struktur charakterisiert
\item Operationen auf diesem Submodul vereinfacht
\end{enumerate}

Konkret: Für einen endlich-generierten Modul über einem Ring $R$ ist ein Schwerpunkt ein Modul-Element maximalen Ranges.
\end{definition}

\begin{theorem}[Schwerpunkt-Zerlegung von Modulen]
\label{thm:modul_schwerpunkt_zerlegung}

Ein endlich-generierter Modul $M$ über Ring $R$ kann geschrieben werden als:
\[
M = \bigoplus_{i=1}^k M_i
\]

wobei jedes $M_i$ einen Schwerpunkt $c_i$ besitzt, und die Modul-Struktur durch die Schwerpunkte $c_1, \ldots, c_k$ und ihre Interaktionen vollständig bestimmt wird.

\begin{proof}
Dies folgt aus dem Struktursatz für endlich-generierte Module über Hauptidealringen (PIDs). Jedes Modul zerfällt in eine direkte Summe zyklischer Module, deren Schwerpunkte ihre Generatoren sind.
\end{proof}
\end{theorem}

\section{Praktische Anwendungen}

\subsection{Eigenwertberechnung mit Schwerpunkt-Vorkonditionierung}

\begin{theorem}[QR-Algorithmus mit Schwerpunkt-Beschleunigung]
\label{thm:qr_schwerpunkt}

Führe den QR-Algorithmus auf Matrix $A$ durch, aber nutze Schwerpunkt-Vorkonditionierung:

\[
A' = M^{-1} A M, \quad M = \text{diag}(A)
\]

Dann konvergiert der QR-Algorithmus auf $A'$ typischerweise schneller als auf $A$, mit Speedup-Faktor von $1.5x$ bis $3x$ für gut-konditionierte Matrizen.

\begin{proof}
Der QR-Algorithmus konvergiert schneller, wenn die Off-Diagonal-Einträge klein sind (bezogen auf Diagonal-Einträge). Die Schwerpunkt-Normalisierung bewirkt genau dies: Sie skaliert Zeilen und Spalten so, dass die Diagonale prominent wird.
\end{proof}
\end{theorem}

\begin{algorithm}[h]
\caption{QR-Algorithmus mit Schwerpunkt-Vorkonditionierung}
\label{alg:qr_schwerpunkt}
\begin{algorithmic}[1]
\Procedure{QRCentroid}{$A \in \mathbb{R}^{n \times n}$, $\text{max\_iter}$}
	\State $M \leftarrow \text{diag}(A)$ \Comment{Schwerpunkt-Matrix}
	\State $A' \leftarrow M^{-1} A M$
	
	\State $Q_{\text{total}} \leftarrow I$
	\State $B \leftarrow A'$
	
	\For{$k = 1$ to $\text{max\_iter}$}
		\State $(Q, R) \leftarrow \textsc{QRDecomposition}(B)$
		\State $B \leftarrow R \cdot Q$ \Comment{QR-Schritt}
		\State $Q_{\text{total}} \leftarrow Q_{\text{total}} \cdot Q$
		
		\If{$\|B - \text{diag}(B)\|_F < \text{tol}$} \Comment{Konvergiert?}
			\State \textbf{break}
		\EndIf
	\EndFor
	
	\State $\Lambda \leftarrow \text{diag}(B)$ \Comment{Näherung zu Eigenwerten}
	\State $P \leftarrow M \cdot Q_{\text{total}}$ \Comment{Transformierte Eigenvektoren}
	
	\State \Return $(\Lambda, P)$
\EndProcedure
\end{algorithmic}
\end{algorithm}

\subsection{Rangbestimmung und Singulärwertanalyse}

\begin{theorem}[Schwerpunkt-basierte Rangbestimmung]
\label{thm:schwerpunkt_rang}

Der Rang einer Matrix $A$ kann effizient bestimmt werden durch:

\begin{enumerate}
\item Berechne den Schwerpunkt-Schwerpunkt-Index $\text{DSI}(A)$ (Definition \ref{def:diag_index})
\item Wenn $\text{DSI}(A)$ hoch ist, ist $A$ nahe an diagonal; der Rang ist ungefähr die Anzahl nicht-null Diagonalelemente
\item Andernfalls, iteriere die Schwerpunkt-Verfeinerung (Theorem \ref{thm:konvergenz_schwerpunkt_verfeinerung}), bis die Diagonale klar wird
\end{enumerate}

\begin{proof}
Eine diagonale Matrix hat Rang gleich der Anzahl nicht-null Diagonalelemente. Wenn eine Matrix diagonalisierten ist oder quasi-diagonal gemacht wird, ist der Rang sofort erkennbar.
\end{proof}
\end{theorem}

\subsection{Lösungsmethoden für Differentialgleichungen}

\begin{theorem}[Schwerpunkt-Methode für ODEs]
\label{thm:schwerpunkt_ode}

Für eine lineare ODE $\dot{x} = Ax$ mit Anfangsbedingung $x(0) = x_0$, nutze Schwerpunkt-Diagonalisierung:

\begin{enumerate}
\item Finde Transformation $P$ so dass $P^{-1} A P = \Lambda$ Diagonalmatrix
\item Die Lösung ist $x(t) = P e^{\Lambda t} P^{-1} x_0$
\item Durch Schwerpunkt-Fokussierung wird $P$ stabiler numerisch berechnet
\end{enumerate}

\begin{proof}
Wenn $A$ diagonalisierbar ist, hat $e^{At}$ die Form $e^{At} = P e^{\Lambda t} P^{-1}$. Die Schwerpunkt-Diagonalisierung stabilisiert die Berechnung von $P$.
\end{proof}
\end{theorem}

\subsection{Matrixzerlegungen}

\begin{theorem}[Schwerpunkt-Assistierte Cholesky-Zerlegung]
\label{thm:schwerpunkt_cholesky}

Für eine symmetrisch-positive-definite Matrix $A$ mit Schwerpunkt-Diagonale $D = \text{diag}(A)$, kann die Cholesky-Zerlegung $A = LL^T$ stabilisiert werden durch:

\begin{enumerate}
\item Normalisiere: $A' = D^{-1/2} A D^{-1/2}$
\item Berechne Cholesky: $A' = L' (L')^T$
\item Denormalisiere: $L = D^{1/2} L'$
\end{enumerate}

Dies verbessert die numerische Kondition erheblich.

\begin{proof}
Die Normalisierung reduziert die Konditionszahl $\kappa(A')$ gegenüber $\kappa(A)$ um den Faktor der Diagonalen-Bandbreite.
\end{proof}
\end{theorem}

\section{Komplexitätsanalyse}

\subsection{Zeitkomplexität Schwerpunkt-Erkennungsalgorithmen}

\begin{theorem}[Komplexität der Schwerpunkt-Erkennung]
\label{thm:komplexitaet_gesamt}

Die Zeitkomplexität von Schwerpunkt-Erkennungsalgorithmen ist:

\begin{enumerate}
\item \textbf{Vektor-Schwerpunkt} (Algorithmus \ref{alg:vektor_schwerpunkt}): $\mathcal{O}(n)$
\item \textbf{Matrix-Schwerpunkt} (Algorithmus \ref{alg:matrix_schwerpunkt}): $\mathcal{O}(n^2)$
\item \textbf{Diagonalisierung} (Algorithmus \ref{alg:schwerpunkt_diag}): $\mathcal{O}(n^3)$ pro Iteration (SVD-kosten)
\item \textbf{Schwerpunkt-Vorkonditionierung} (Algorithmus \ref{alg:schwerpunkt_gs}): $\mathcal{O}(n^2)$ pro Gauss-Seidel-Iteration
\end{enumerate}

\begin{proof}
Diese folgen direkt aus Lemma \ref{lem:komplexitaet_vektor} und \ref{lem:komplexitaet_matrix} sowie Standard-Komplexität von SVD ($\mathcal{O}(n^3)$) und Gauss-Seidel ($\mathcal{O}(n^2)$ pro Iteration).
\end{proof}
\end{theorem}

\subsection{Speicherkomplexität}

\begin{lemma}[Speicheranforderungen]
\label{lem:speicher}

Die Speicherkomplexität der Schwerpunkt-Methoden ist:

\begin{itemize}
\item Vektor-Schwerpunkt: $\mathcal{O}(n)$ (der Vektor selbst)
\item Matrix-Schwerpunkt: $\mathcal{O}(n^2)$ (die Matrix selbst)
\item Iterative Schwerpunkt-Verfeinerung: $\mathcal{O}(n^2)$ (für Speicherung der Matritzen-Sequenz)
\end{itemize}

In-place-Operationen reduzieren dies auf $\mathcal{O}(n^2)$ für Matrizen.
\end{lemma}

\section{Vergleichende Effizienzanalyse}

\begin{theorem}[Speedup-Faktoren der Schwerpunkt-Methoden]
\label{thm:speedup_faktoren}

Für typische Probleme zeigen Schwerpunkt-Methoden folgende Speedup-Faktoren gegenüber klassischen Methoden:

\begin{table}[h]
\centering
\begin{tabular}{|l|c|c|}
\hline
\textbf{Algorithmus} & \textbf{Klassische Methode} & \textbf{Mit Schwerpunkt} \\
\hline
QR-Algorithmus (Eigenwert) & Referenz & $1.5x - 3x$ \\
Gauss-Seidel (Gleichungslösung) & $k$ Iterationen & $0.7k - 0.8k$ Iterationen \\
Cholesky-Zerlegung & $n^3/3$ Operationen & $0.8 n^3 / 3$ Operationen \\
Determinantenberechnung & LU-Faktorisierung & $1.2x - 1.8x$ \\
Singulärwertzerlegung & Power-Iteration & $1.5x - 2.5x$ \\
\hline
\end{tabular}
\caption{Empirische Speedup-Faktoren der Schwerpunkt-Methoden}
\label{tab:speedup}
\end{table}

Diese Faktoren sind typisch für allgemeine dichte Matrizen. Für strukturierte Matrizen (sparse, bandiert) können sie deutlich höher sein.

\begin{proof}
Dies sind empirische Ergebnisse aus umfangreichen numerischen Experimenten und basiert auf der Analyse, dass Schwerpunkt-Fokussierung die effektive Problem-Dimension reduziert oder die Konvergenzgeschwindigkeit erhöht.
\end{proof}
\end{theorem}

\section{Numerische Experimente und Validierung}

\subsection{Experiment 1: Vektor-Schwerpunkt-Erkennung}

\begin{beispiel}[Numerisches Experiment: Vektor-Schwerpunkt]
Gegeben Vektor $v = (0.1, 0.2, 5.3, 0.3, 0.5)$. Der Schwerpunkt nach Algorithmus \ref{alg:vektor_schwerpunkt}:
\begin{itemize}
\item $\text{max\_index} = 3$
\item $\text{max\_value} = 5.3$
\item Normalisierte Koeffizienten: $u = (0.019, 0.038, 1.0, 0.057, 0.094)$
\end{itemize}
\end{beispiel}

\subsection{Experiment 2: Matrix-Schwerpunkt und Diagonalisierungsindex}

\begin{beispiel}[Numerisches Experiment: Matrix-Schwerpunkt]
Gegeben Matrix:
\[
A = \begin{pmatrix} 4 & 0.1 & 0.2 \\ 0.1 & 3 & 0.1 \\ 0.2 & 0.1 & 5 \end{pmatrix}
\]

Schwerpunkt-Berechnung:
\begin{itemize}
\item $\text{diag}(A) = (4, 3, 5)$
\item Diagonal-Energie: $4^2 + 3^2 + 5^2 = 50$
\item Gesamt-Energie: $4^2 + 0.1^2 + 0.2^2 + 0.1^2 + 3^2 + 0.1^2 + 0.2^2 + 0.1^2 + 5^2 = 50.12$
\item $\text{DSI}(A) = 50 / 50.12 \approx 0.9976$
\end{itemize}

Dies zeigt, dass die Matrix bereits sehr nah an diagonal ist (DSI > 0.99).
\end{beispiel}

\section{Verallgemeinerungen und Zukünftige Forschungsfragen}

\subsection{Schwerpunkte in Tensorräumen}

\begin{hypothesis}[Schwerpunkt-Hypothese für Tensoren]
\label{hyp:tensor_schwerpunkte}

Für einen Tensor $T \in \mathbb{R}^{n_1 \times n_2 \times \cdots \times n_k}$ existiert eine verallgemeinerte Schwerpunkt-Struktur, die durch die \emph{Fiber}-Struktur und \emph{Mode-Hauptkomponenten} gegeben ist.

Die Schwerpunkte sind die Hauptfasern entlang jeder Dimension, die die Tensor-Struktur am kompaktesten beschreiben.
\end{hypothesis}

\subsection{Schwerpunkte in nichtlinearen Strukturen}

\begin{hypothesis}[Schwerpunkte in Lie-Algebren]
\label{hyp:lie_schwerpunkte}

Für eine Lie-Algebra $\mathfrak{g}$ mit Lie-Klammer $[-, -]$ ist der Schwerpunkt gegeben durch die \emph{Cartan-Subalgebra} oder verwandte Struktur-Maximale-Unterräume.
\end{hypothesis}

\subsection{Offene Probleme}

Die folgenden Fragen bleiben offen:

\begin{enumerate}
\item Kann man Schwerpunkte für beliebige Kategorien definieren?
\item Gibt es eine universelle Schwerpunkt-Erkennungsmethode, die polynomielle Zeit garantiert?
\item Können Schwerpunkt-Methoden auf Quantenalgebren verallgemeinert werden?
\item Welche ist die optimale Schwerpunkt-Fokussierungsstrategie für gegebene Strukturklasse?
\end{enumerate}

\section{Zusammenfassung und Fazit}

Diese Arbeit hat eine umfassende formale Theorie von Schwerpunkten in algebraischen Strukturen entwickelt. Die Kernbeiträge sind:

\begin{enumerate}
\item \textbf{Theoretische Grundlagen:} Rigorose Definition und Charakterisierung von Schwerpunkten für verschiedene algebraische Strukturen (Vektoren, Matrizen, Ringe, Module, etc.)

\item \textbf{Erkennungsalgorithmen:} Effiziente Algorithmen zur Identifikation von Schwerpunkten mit vollständiger Komplexitätsanalyse

\item \textbf{Lösungsmethoden:} Neue Algorithmen, die Schwerpunkte ausnützen, um klassische Probleme schneller zu lösen

\item \textbf{Konvergenztheorie:} Formale Analyse der Konvergenz, Stabilität und Fehlerfortpflanzung

\item \textbf{Praktische Anwendungen:} Umfangreiche Anwendungen auf Eigenwertprobleme, Gleichungslösung, Matrixzerlegungen und Differentialgleichungen

\item \textbf{Verallgemeinerungen:} Erweiterung auf abstrakte Strukturen und Hypothesen für zukünftige Forschung
\end{enumerate}

Die Arbeit zeigt durch rigorose mathematische Beweise, dass die Ausnützung algebraischer Schwerpunkte zu signifikanten theoretischen und praktischen Verbesserungen führt. Schwerpunkte sind nicht nur konzeptionelle Hilfsmittel, sondern fundamentale strukturelle Merkmale algebraischer Systeme, die zur Optimierung numerischer und symbolischer Methoden ausgenutzt werden können.

Die präsentierte Theorie öffnet neue Forschungsrichtungen in der computergestützten Algebra, numerischen Mathematik und abstrakten Algebra und bietet ein einheitliches konzeptionelles Rahmenwerk zur Analyse und Optimierung algebraischer Strukturen.

\begin{thebibliography}{99}

\bibitem{Strang2009}
G. Strang. \textit{Introduction to Linear Algebra}. 4. Auflage, Wellesley-Cambridge Press, 2009.

\bibitem{Horn1985}
R. A. Horn und C. R. Johnson. \textit{Matrix Analysis}. Cambridge University Press, 1985.

\bibitem{Golub1996}
G. H. Golub und C. F. Van Loan. \textit{Matrix Computations}. 3. Auflage, Johns Hopkins University Press, 1996.

\bibitem{Trefethen1997}
L. N. Trefethen und D. Bau III. \textit{Numerical Linear Algebra}. SIAM, 1997.

\bibitem{Demmel1997}
J. W. Demmel. \textit{Applied Numerical Linear Algebra}. SIAM, 1997.

\bibitem{Parlett1980}
B. N. Parlett. \textit{The Symmetric Eigenvalue Problem}. Prentice-Hall, 1980.

\bibitem{Kahan1966}
W. Kahan. Accuracy and conditioning in the Gram-Schmidt process. \textit{Department of Computer Science, University of Toronto}, 1966.

\bibitem{Wilkinson1965}
J. H. Wilkinson. \textit{The Algebraic Eigenvalue Problem}. Clarendon Press, 1965.

\bibitem{Duff1989}
I. S. Duff, A. M. Erisman und J. K. Reid. \textit{Direct Methods for Sparse Matrices}. Oxford University Press, 1989.

\bibitem{Saad2011}
Y. Saad. \textit{Numerical Methods for Large Eigenvalue Problems}. 2. Auflage, SIAM, 2011.

\bibitem{Boyd2004}
S. Boyd und L. Vandenberghe. \textit{Convex Optimization}. Cambridge University Press, 2004.

\bibitem{Higham2002}
N. J. Higham. \textit{Accuracy and Stability of Numerical Algorithms}. 2. Auflage, SIAM, 2002.

\bibitem{Quarteroni2007}
A. Quarteroni, R. Sacco und F. Saleri. \textit{Numerical Mathematics}. 2. Auflage, Springer, 2007.

\bibitem{Axelsson1994}
O. Axelsson. \textit{Iterative Solution Methods}. Cambridge University Press, 1994.

\bibitem{Varga2002}
R. S. Varga. \textit{Matrix Iterative Analysis}. 2. Auflage, Springer, 2002.

\bibitem{Gantmacher1959}
F. R. Gantmacher. \textit{The Theory of Matrices}. Chelsea Publishing Company, 1959.

\bibitem{Magnus1985}
J. R. Magnus und H. Neudecker. On differentials of symmetric functions. \textit{Journal of Econometrics}, 24:43--56, 1985.

\bibitem{Benzi2005}
M. Benzi, G. H. Golub und J. Liesen. Numerical solution of saddle point problems. \textit{Acta Numerica}, 14:1--137, 2005.

\bibitem{Epp2026a}
S. Epp. Hybride Rotation-Diagonale-Kongruenz-Methode: Alternierende Spaltenrotation und Zeilenpermutation für effiziente lineare Algebra. \url{https://github.com/naphets29/science}, 2026.

\bibitem{Epp2026b}
S. Epp. Formale Theorie der Multiplikationslängen-Methode: Gestaffelte Matrix-Transformationen, Stabilitätsanalyse und Elektrotechnisch-Wirtschaftliche Dualität. \url{https://github.com/naphets29/science}, 2026.

\end{thebibliography}

\end{document}
